% GENERATED FILE. Edit canonical lessons, then run npm run generate:books.
% canonical-source-hash: fnv1a64:b25ae196711fa247
% canonical-lessons: MR-C484-rog, MR-C484-samsarga, MR-C484-jantu, MR-C484-visanu, MR-C484-sath

\chapter{Germs and Infections}
\label{ch:mr-germs-and-infections}
\hlchaptermodality{484}

\begin{chapteropening}
\textbf{By the end of this chapter:} \emph{I can talk about a disease, an infection, germs and an epidemic.}

Complete the last lesson of chapter 484: I can talk about a disease, an infection, germs and an epidemic.
\end{chapteropening}

\section[\texorpdfstring{\mr{रोग} (rog)}{रोग (rog)}]{\mr{रोग} (rog) — a disease}

\label{lesson:MR-C484-rog}



\begin{quote}
Before the new one: say the Marathi for to vomit, then the Marathi for to be digested.
\end{quote}

\subsection*{You'll want to know: \mr{रोग}}
\textbf{\mr{रोग}} — \emph{rog} — \textquotedblleft{}a disease\textquotedblright{}.

\begin{quote}
\mr{स्वच्छ} \mr{पाण्यामुळे} \mr{बरेच} \mr{रोग} \mr{टळतात}.
\end{quote}

\begin{quote}
\emph{svaccha pāṇyāmuḷe barec rog ṭaḷtāt.}
\end{quote}

\begin{quote}
Clean water keeps many diseases away.
\end{quote}

Masculine.

Masculine; plural \mr{रोग}.

\begin{cousinweb}
\textbf{\mr{रोग}} (\emph{rog}) is a tatsama, taken from Sanskrit \emph{roga}, \textquotedblleft{}disease\textquotedblright{} (CDIAL 10829).
\end{cousinweb}

\begin{grammarlens}[title={What you've built}]
One more word for describing what you lived and felt.
\end{grammarlens}

\subsection*{Your turn}
\begin{itemize}
\raggedright
  \item \emph{Say it:} \emph{rog}
  \item \emph{Say it:} \emph{rog}, once more
  \item \emph{Say it:} \emph{svaccha pāṇyāmuḷe barec rog ṭaḷtāt.} — Clean water keeps many diseases away.
  \item \emph{Recall:} say the Marathi for to vomit, then the Marathi for to be digested, then say \emph{rog} again
\end{itemize}

\begin{tcolorbox}[breakable,colback=teal!4,colframe=teal!35!black,title={Before you move on}]
What does \mr{रोग} mean? (\textquotedblleft{}A disease\textquotedblright{}.) Say it once more.
\end{tcolorbox}
\section[\texorpdfstring{\mr{संसर्ग} (saṁsarga)}{संसर्ग (saṁsarga)}]{\mr{संसर्ग} (saṁsarga) — an infection; contagion}

\label{lesson:MR-C484-samsarga}



\begin{quote}
Before the new one: say the Marathi for to be digested, then the Marathi for a disease.
\end{quote}

\subsection*{You'll want to know: \mr{संसर्ग}}
\textbf{\mr{संसर्ग}} — \emph{saṁsarga} — \textquotedblleft{}an infection; contagion\textquotedblright{}.

\begin{quote}
\mr{जखमेत} \mr{संसर्ग} \mr{होऊ} \mr{नये} \mr{म्हणून} \mr{मलम} \mr{लावलं}.
\end{quote}

\begin{quote}
\emph{jakhmet saṁsarga hoū naye mhaṇūn malam lāvlaṁ.}
\end{quote}

\begin{quote}
We put ointment on the wound so that it would not get infected.
\end{quote}

Masculine.

\mr{संसर्ग} \mr{होणे}, to get infected; \mr{संसर्गजन्य}, contagious.

\begin{cousinweb}
\textbf{\mr{संसर्ग}} (\emph{saṁsarga}) is a tatsama, taken from Sanskrit \emph{saṁsarga}, \textquotedblleft{}mixture, intercourse\textquotedblright{}, that is, contact (CDIAL 13033).
\end{cousinweb}

\begin{grammarlens}[title={What you've built}]
One more word for describing what you lived and felt.
\end{grammarlens}

\subsection*{Your turn}
\begin{itemize}
\raggedright
  \item \emph{Say it:} \emph{saṁsarga}
  \item \emph{Say it:} \emph{saṁsarga}, once more
  \item \emph{Say it:} \emph{jakhmet saṁsarga hoū naye mhaṇūn malam lāvlaṁ.} — We put ointment on the wound so that it would not get infected.
  \item \emph{Recall:} say the Marathi for to be digested, then the Marathi for a disease, then say \emph{saṁsarga} again
\end{itemize}

\begin{tcolorbox}[breakable,colback=teal!4,colframe=teal!35!black,title={Before you move on}]
What does \mr{संसर्ग} mean? (\textquotedblleft{}An infection; contagion\textquotedblright{}.) Say it once more.
\end{tcolorbox}
\section[\texorpdfstring{\mr{जंतू} (jantū)}{जंतू (jantū)}]{\mr{जंतू} (jantū) — a germ}

\label{lesson:MR-C484-jantu}



\begin{quote}
Before the new one: say the Marathi for a disease, then the Marathi for an infection.
\end{quote}

\subsection*{You'll want to know: \mr{जंतू}}
\textbf{\mr{जंतू}} — \emph{jantū} — \textquotedblleft{}a germ\textquotedblright{}.

\begin{quote}
\mr{साबणाने} \mr{हात} \mr{धुतले} \mr{की} \mr{जंतू} \mr{जातात}.
\end{quote}

\begin{quote}
\emph{sābṇāne hāt dhutle kī jantū jātāt.}
\end{quote}

\begin{quote}
Washing your hands with soap gets rid of germs.
\end{quote}

Masculine.

Masculine; plural \mr{जंतू}. \mr{जंतुनाशक}, a disinfectant.

\begin{cousinweb}
\textbf{\mr{जंतू}} (\emph{jantū}) is a tatsama, taken from Sanskrit \emph{jantu}, \textquotedblleft{}creature; any animal of the lowest order\textquotedblright{} (CDIAL 5110).
\end{cousinweb}

\begin{grammarlens}[title={What you've built}]
One more word for describing what you lived and felt.
\end{grammarlens}

\subsection*{Your turn}
\begin{itemize}
\raggedright
  \item \emph{Say it:} \emph{jantū}
  \item \emph{Say it:} \emph{jantū}, once more
  \item \emph{Say it:} \emph{sābṇāne hāt dhutle kī jantū jātāt.} — Washing your hands with soap gets rid of germs.
  \item \emph{Recall:} say the Marathi for a disease, then the Marathi for an infection, then say \emph{jantū} again
\end{itemize}

\begin{tcolorbox}[breakable,colback=teal!4,colframe=teal!35!black,title={Before you move on}]
What does \mr{जंतू} mean? (\textquotedblleft{}A germ\textquotedblright{}.) Say it once more.
\end{tcolorbox}
\section[\texorpdfstring{\mr{विषाणू} (viṣāṇū)}{विषाणू (viṣāṇū)}]{\mr{विषाणू} (viṣāṇū) — a virus}

\label{lesson:MR-C484-visanu}



\begin{quote}
Before the new one: say the Marathi for an infection, then the Marathi for a germ.
\end{quote}

\subsection*{You'll want to know: \mr{विषाणू}}
\textbf{\mr{विषाणू}} — \emph{viṣāṇū} — \textquotedblleft{}a virus\textquotedblright{}.

\begin{quote}
\mr{सर्दी} \mr{एका} \mr{विषाणूमुळे} \mr{होते}.
\end{quote}

\begin{quote}
\emph{sardī ekā viṣāṇūmuḷe hote.}
\end{quote}

\begin{quote}
A cold is caused by a virus.
\end{quote}

Masculine.

\mr{विष}, poison, + \mr{अणू}, a tiny particle.

\begin{cousinweb}
\textbf{\mr{विषाणू}} (\emph{viṣāṇū}) is a modern learned compound built from two Sanskrit words, \emph{viṣa}, \textquotedblleft{}poison\textquotedblright{}, and \emph{aṇu}, \textquotedblleft{}minute, atomic\textquotedblright{} (Monier-Williams).
\end{cousinweb}

\begin{grammarlens}[title={What you've built}]
One more word for describing what you lived and felt.
\end{grammarlens}

\subsection*{Your turn}
\begin{itemize}
\raggedright
  \item \emph{Say it:} \emph{viṣāṇū}
  \item \emph{Say it:} \emph{viṣāṇū}, once more
  \item \emph{Say it:} \emph{sardī ekā viṣāṇūmuḷe hote.} — A cold is caused by a virus.
  \item \emph{Recall:} say the Marathi for an infection, then the Marathi for a germ, then say \emph{viṣāṇū} again
\end{itemize}

\begin{tcolorbox}[breakable,colback=teal!4,colframe=teal!35!black,title={Before you move on}]
What does \mr{विषाणू} mean? (\textquotedblleft{}A virus\textquotedblright{}.) Say it once more.
\end{tcolorbox}
\section[\texorpdfstring{\mr{साथ} (sāth)}{साथ (sāth)}]{\mr{साथ} (sāth) — an epidemic}

\label{lesson:MR-C484-sath}



\begin{quote}
Before the new one: say the Marathi for a germ, then the Marathi for a virus.
\end{quote}

\subsection*{You'll want to know: \mr{साथ}}
\textbf{\mr{साथ}} — \emph{sāth} — \textquotedblleft{}an epidemic\textquotedblright{}.

\begin{quote}
\mr{पावसाळ्यात} \mr{गावात} \mr{तापाची} \mr{साथ} \mr{आली}.
\end{quote}

\begin{quote}
\emph{pāvsāḷyāt gāvāt tāpācī sāth ālī.}
\end{quote}

\begin{quote}
In the monsoon a fever epidemic came to the village.
\end{quote}

Feminine.

\mr{साथीचा} \mr{रोग}, an infectious disease. A second \mr{साथ} means company: \mr{साथ} \mr{देणे}, to stand by someone.

\begin{cousinweb}
\textbf{\mr{साथ}} (\emph{sāth}) is inherited in the sense \textquotedblleft{}company\textquotedblright{}: CDIAL lists Marathi \emph{sāth}, \textquotedblleft{}companionship\textquotedblright{}, under Sanskrit \emph{sārtha}, \textquotedblleft{}caravan, company\textquotedblright{} (CDIAL 13364). Whether the epidemic sense is the same word is not settled.
\end{cousinweb}

\begin{grammarlens}[title={What you've built}]
One more word for describing what you lived and felt.
\end{grammarlens}

\subsection*{Your turn}
\begin{itemize}
\raggedright
  \item \emph{Say it:} \emph{sāth}
  \item \emph{Say it:} \emph{sāth}, once more
  \item \emph{Say it:} \emph{pāvsāḷyāt gāvāt tāpācī sāth ālī.} — In the monsoon a fever epidemic came to the village.
  \item \emph{Recall:} say the Marathi for a germ, then the Marathi for a virus, then say \emph{sāth} again
\end{itemize}

\begin{tcolorbox}[breakable,colback=teal!4,colframe=teal!35!black,title={Before you move on}]
What does \mr{साथ} mean? (\textquotedblleft{}An epidemic\textquotedblright{}.) Say it once more.
\end{tcolorbox}
